\documentclass[10pt,a4paper]{amsart}
\usepackage[margin=19mm]{geometry}
\usepackage{amsmath,amssymb,mathtools}
\usepackage[scr=rsfs,cal=euler]{mathalfa}
\usepackage{xfrac}
\usepackage[unicode,hidelinks,bookmarks=false]{hyperref}
\newcommand{\ie}{i.e.\ }
\newcommand{\cf}{cf.\ }
\newcommand{\resp}{respectively\ }
\newcommand{\Subsection}{\subsection}
\providecommand{\nobreakdash}{\nobreak\hbox{-}}
\setlength{\emergencystretch}{2em}
\multlinegap0pt
\usepackage{lmodern}
\usepackage{mathtools}
\usepackage{mathrsfs}
\usepackage{enumitem}
\usepackage{booktabs}
\usepackage{longtable}
\usepackage{array}
\usepackage{needspace}
\usepackage{cleveref}
\usepackage{tikz}
\usetikzlibrary{matrix,positioning,calc,arrows.meta,fit,shapes.geometric,backgrounds}
\setlength{\emergencystretch}{2em}
\allowdisplaybreaks[2]
\numberwithin{equation}{section}
\hypersetup{
  pdftitle={Closure Atlases and Local-to-Global Obstructions in Finite Closure Systems},
  pdfauthor={Jaehwan Kim}
}
\newtheorem{theorem}{Theorem}[section]
\newtheorem{lemma}[theorem]{Lemma}
\newtheorem{proposition}[theorem]{Proposition}
\newtheorem{corollary}[theorem]{Corollary}
\newtheorem{definition}[theorem]{Definition}
\newtheorem{example}[theorem]{Example}
\newtheorem{counterexample}[theorem]{Counterexample}
\newtheorem{construction}[theorem]{Construction}
\newtheorem{remark}[theorem]{Remark}
\newcommand{\Pow}{\mathcal P}
\newcommand{\Cl}{\operatorname{cl}}
\newcommand{\Fix}{\operatorname{Fix}}
\newcommand{\C}{\mathcal C}
\newcommand{\X}{\mathfrak X}
\newcommand{\Y}{\mathfrak Y}
\newcommand{\F}{\mathfrak F}
\newcommand{\A}{\mathfrak A}
\newcommand{\Sig}{\Sigma}
\newcommand{\Obs}{\operatorname{Obs}}
\newcommand{\CA}{C_{\mathfrak A}}
\newcommand{\GA}{G_{\mathfrak A}}
\newcommand{\Tcup}{T^{\cup}}
\newcommand{\restr}[1]{\upharpoonright_{#1}}
\begin{document}
\title[Closure Atlases and Local-to-Global Obstructions in Finite C]{Closure Atlases and Local-to-Global Obstructions in Finite Closure Systems}
\author{Jaehwan Kim}
\date{}
\maketitle
\noindent\emph{Source and licence.} Closure Atlases and Local-to-Global Obstructions in Finite Closure Systems. Version arXiv:2606.24909v1 (CC BY 4.0). This material is distributed under the Creative Commons Attribution 4.0 International licence, http://creativecommons.org/licenses/by/4.0/, as recorded in the arXiv OAI licence metadata for this exact version and shown on the abstract page https://arxiv.org/abs/2606.24909v1. Full rights notice: licenses/arxiv-2606.24909-CC-BY-4.0.txt.\par\smallskip
\noindent\textit{Editorial scope.} Counted unit: the original \texttt{theorem} environment \texttt{thm:full-truth-region-representation} at source lines 343--363 together with its complete proof at lines 365--383; the delivered document keeps source lines 238--384 so that every referenced label and every citation resolves inside it. Native editable source is the paper's own concatenated TeX; the AMS wrapper is editorial formatting only. No publisher class, upstream script or unrelated artwork is redistributed.
\section{Finite Closure Systems and Indexed Truth Spaces}
\label{sec:finite-closure-systems}

We begin with the single-system case.  The finiteness assumption is not needed for every definition in this section, but it will be used later when the obstruction theorem is converted into an explicit finite algorithm.  The section has two roles: first, to fix the closure consequence relation of a finite closure operator; second, to compare that relation with the weaker information visible through a chosen indexed family of closed sets, called closed theories in the intended logical reading.  This distinction is the base case for every later comparison between local and global closure data.

\begin{definition}[Finite closure system, consequence, and closed theories]
A \emph{finite closure system} is a pair
\[
  \C=(\Sig,\Cl),
\]
where \(\Sig\) is a finite set of sentences and
\[
  \Cl:\Pow(\Sig)\to\Pow(\Sig)
\]
is extensive, monotone, and idempotent; that is, for all \(S,U\subseteq\Sig\),
\begin{enumerate}[label=(\roman*)]
  \item \(S\subseteq\Cl(S)\),
  \item if \(S\subseteq U\), then \(\Cl(S)\subseteq\Cl(U)\),
  \item \(\Cl(\Cl(S))=\Cl(S)\).
\end{enumerate}
For \(S\subseteq\Sig\) and \(\varphi\in\Sig\), write
\[
  S\vdash_{\C}\varphi
  \quad\Longleftrightarrow\quad
  \varphi\in\Cl(S).
\]
Equivalently, one may write \(S\vdash_{\Cl}\varphi\).  When the closure system is clear from context, we simply write \(S\vdash\varphi\).

A subset \(T\subseteq\Sig\) is a \emph{closed theory} of \(\C\) if \(\Cl(T)=T\).  Write
\[
  \Fix(\Cl)=\{T\subseteq\Sig: \Cl(T)=T\}
\]
for the set of all closed theories of \(\C\).
\end{definition}

\begin{remark}[Raw sets and closed theories]
A subset \(S\subseteq\Sig\) is an arbitrary set of sentences; it need not be closed. A closed theory is a stabilized object \(T=\Cl(T)\). Thus \(S\vdash\varphi\) does not mean that \(\varphi\in S\). It means that \(\varphi\) belongs to the closed theory generated by \(S\).
\end{remark}

\begin{definition}[Finite indexed truth space and truth regions]
Let \(\C=(\Sig,\Cl)\) be a finite closure system. A \emph{finite indexed truth space} over \(\C\) is a tuple
\[
  \X=(\Sig,\Cl,I,\tau),
\]
where \(I\) is a finite index set and
\[
  \tau:I\to\Fix(\Cl)
\]
assigns a closed theory to each index.  The elements of \(I\) are called \emph{indices} or, when helpful, \emph{contexts}. In the main examples and representation theorems below, the relevant index sets are nonempty.

For \(\varphi\in\Sig\), the \emph{truth region} of \(\varphi\) in \(\X\) is
\[
  V_{\X}(\varphi)=\{i\in I:\varphi\in\tau(i)\}.
\]
For \(S\subseteq\Sig\), define
\[
  V_{\X}(S)=\{i\in I:S\subseteq\tau(i)\}
  =\bigcap_{\psi\in S}V_{\X}(\psi),
\]
with \(V_{\X}(\varnothing)=I\).  When \(\X\) is clear from context, we may write \(V(\varphi)\) and \(V(S)\).
\end{definition}

\begin{definition}[Region consequence]
For \(S\subseteq\Sig\) and \(\varphi\in\Sig\), write
\[
  S\models_{\X}\varphi
\]
if
\[
  V_{\X}(S)\subseteq V_{\X}(\varphi).
\]
Thus every selected index whose assigned closed theory contains all sentences in \(S\) also contains \(\varphi\). This relation depends on the chosen indexed truth space \(\X\). It should not be confused with closure consequence \(S\vdash_{\C}\varphi\).
\end{definition}

\begin{lemma}[Truth-region soundness]
\label{lem:truth-region-soundness}
Let \(\X=(\Sig,\Cl,I,\tau)\) be a finite indexed truth space over \(\C=(\Sig,\Cl)\). For all \(S\subseteq\Sig\) and \(\varphi\in\Sig\),
\[
  S\vdash_{\C}\varphi
  \quad\Longrightarrow\quad
  V_{\X}(S)\subseteq V_{\X}(\varphi).
\]
Equivalently, \(S\vdash_{\C}\varphi\) implies \(S\models_{\X}\varphi\).
\end{lemma}

\begin{proof}
Assume \(S\vdash_{\C}\varphi\), so \(\varphi\in\Cl(S)\). Let \(i\in V_{\X}(S)\). Then \(S\subseteq\tau(i)\). By monotonicity of \(\Cl\),
\[
  \Cl(S)\subseteq\Cl(\tau(i)).
\]
Since \(\tau(i)\) is a closed theory, \(\Cl(\tau(i))=\tau(i)\). Hence \(\Cl(S)\subseteq\tau(i)\). Since \(\varphi\in\Cl(S)\), it follows that \(\varphi\in\tau(i)\), and therefore \(i\in V_{\X}(\varphi)\). This proves the inclusion.
\end{proof}

\begin{remark}[Soundness need not be completeness]
Lemma~\ref{lem:truth-region-soundness} says that closure consequence is always sound for truth-region inclusion. The converse need not hold for an arbitrary indexed truth space: the chosen index set may omit closed theories that distinguish \(S\) from \(\varphi\). The exact converse is recovered when all closed theories are used as indices.
\end{remark}

\begin{definition}[Full indexed truth space]
Let \(\C=(\Sig,\Cl)\) be a finite closure system. The \emph{full indexed truth space} of \(\C\) is
\[
  \F_{\C}=(\Sig,\Cl,\Fix(\Cl),\mathrm{id}),
\]
where the index set is \(\Fix(\Cl)\) and the assignment map is the identity, \(\mathrm{id}(T)=T\). Thus every closed theory is used as an index. Since \(\Sig\) is finite, \(\Fix(\Cl)\) is finite; it is nonempty because \(\Cl(\varnothing)\) is closed.
\end{definition}


\begin{theorem}[Full truth-region representation]
\label{thm:full-truth-region-representation}
Let \(\C=(\Sig,\Cl)\) be a finite closure system, and let \(\F_{\C}\) be its full indexed truth space. For every \(S\subseteq\Sig\) and every \(\varphi\in\Sig\),
\[
  V_{\F_{\C}}(S)\subseteq V_{\F_{\C}}(\varphi)
  \quad\Longleftrightarrow\quad
  S\vdash_{\C}\varphi.
\]
Equivalently,
\[
  S\models_{\F_{\C}}\varphi
  \quad\Longleftrightarrow\quad
  S\vdash_{\C}\varphi.
\]
More generally, for \(S,U\subseteq\Sig\),
\[
  V_{\F_{\C}}(S)\subseteq V_{\F_{\C}}(U)
  \quad\Longleftrightarrow\quad
  U\subseteq\Cl(S).
\]
\end{theorem}

\begin{proof}
We first prove the sentence case. If \(S\vdash_{\C}\varphi\), then Lemma~\ref{lem:truth-region-soundness}, applied to the indexed truth space \(\F_{\C}\), gives
\[
  V_{\F_{\C}}(S)\subseteq V_{\F_{\C}}(\varphi).
\]
Conversely, assume
\[
  V_{\F_{\C}}(S)\subseteq V_{\F_{\C}}(\varphi).
\]
By idempotence, \(\Cl(S)\) is closed, so \(\Cl(S)\in\Fix(\Cl)\). By extensiveness, \(S\subseteq\Cl(S)\), hence \(\Cl(S)\in V_{\F_{\C}}(S)\). The assumed inclusion gives \(\Cl(S)\in V_{\F_{\C}}(\varphi)\). Since the assignment in \(\F_{\C}\) is the identity, this means \(\varphi\in\Cl(S)\), i.e. \(S\vdash_{\C}\varphi\).

For the general statement, first assume \(V_{\F_{\C}}(S)\subseteq V_{\F_{\C}}(U)\). For each \(\psi\in U\), we have \(V_{\F_{\C}}(U)\subseteq V_{\F_{\C}}(\psi)\), hence \(V_{\F_{\C}}(S)\subseteq V_{\F_{\C}}(\psi)\). By the sentence case, \(\psi\in\Cl(S)\). Therefore \(U\subseteq\Cl(S)\).

Conversely, assume \(U\subseteq\Cl(S)\), and let \(T\in V_{\F_{\C}}(S)\). Then \(T\in\Fix(\Cl)\) and \(S\subseteq T\). By monotonicity,
\[
  \Cl(S)\subseteq\Cl(T)=T.
\]
Since \(U\subseteq\Cl(S)\), it follows that \(U\subseteq T\), so \(T\in V_{\F_{\C}}(U)\). Therefore \(V_{\F_{\C}}(S)\subseteq V_{\F_{\C}}(U)\).
\end{proof}


\end{document}
